% =============================================================================
% TERMO DE CONSENTIMENTO LIVRE E ESCLARECIDO (TCLE)
% Versão para Agricultores Quilombolas
% =============================================================================

\documentclass[12pt,a4paper,brazil]{article}

\usepackage[utf8]{inputenc}
\usepackage[T1]{fontenc}
\usepackage[brazil]{babel}
\usepackage{geometry}
\usepackage{setspace}
\usepackage{indentfirst}
\usepackage{graphicx}
\usepackage{fancyhdr}

\geometry{a4paper, left=3cm, right=2cm, top=3cm, bottom=2cm}
\onehalfspacing

\pagestyle{fancy}
\fancyhf{}
\rhead{\small TCLE --- Agricultores Quilombolas}
\lhead{\small PPGPI/UFS}
\cfoot{\thepage}

\begin{document}

\begin{center}
{\large\bfseries UNIVERSIDADE FEDERAL DE SERGIPE}\\[0.3cm]
{\normalsize Programa de Pós-Graduação em Ciência da Propriedade Intelectual (PPGPI)}\\[1cm]
{\Large\bfseries TERMO DE CONSENTIMENTO LIVRE E ESCLARECIDO (TCLE)}\\[0.3cm]
{\normalsize Versão para Agricultores Quilombolas}\\[0.5cm]
\noindent\rule{\textwidth}{0.4pt}
\end{center}

\vspace{0.5cm}

\noindent\textbf{Título da Pesquisa.} Modelo Computacional para Salvaguarda de Saberes Agrícolas Tradicionais como Ativos Intangíveis em Comunidades Quilombolas do Semiárido Nordeste~II.

\noindent\textbf{Pesquisador Responsável.} Prof. Dr. Luiz Diego Vidal Santos (UEFS/PPGPI-UFS).

\noindent\textbf{Doutoranda.} Catuxe Varjão de Santana Oliveira.

\noindent\textbf{Orientador.} Prof. Dr. Paulo Roberto Gagliardi (PPGPI/UFS).

\noindent\textbf{Instituição.} Universidade Federal de Sergipe (UFS), Programa de Pós-Graduação em Ciência da Propriedade Intelectual.

\vspace{0.5cm}

\section*{Convite à Participação}

Você está sendo convidado(a) a participar de uma pesquisa sobre os saberes e práticas agrícolas tradicionais da sua comunidade. Este documento traz todas as informações necessárias para que você decida, de forma livre e bem informada, se deseja ou não participar. Leia com atenção (ou ouça a leitura feita pela equipe de pesquisa) e sinta-se à vontade para fazer qualquer pergunta antes de tomar sua decisão.

\section*{Qual é o objetivo da pesquisa?}

A pesquisa tem como objetivo desenvolver e validar uma ferramenta computacional que possa auxiliar as comunidades quilombolas do Território de Identidade Semiárido Nordeste~II (Bahia) na proteção dos seus próprios saberes agrícolas tradicionais. Para isso, será adaptado um questionário internacional sobre práticas de manejo da terra para a linguagem e a realidade das comunidades, com a participação direta de agricultores e mestres de saberes. Com base nesse questionário adaptado, será criado um índice (chamado Índice de Salvaguarda Biocultural) que ajuda a identificar quais saberes podem precisar de atenção mais urgente. A ferramenta ficará à disposição das comunidades e poderá ser utilizada em processos junto ao IPHAN, INPI e outros órgãos, fortalecendo a capacidade das comunidades de proteger e valorizar seus conhecimentos por conta própria.

\section*{O que será pedido a você?}

Dependendo do componente da pesquisa em que você participar, poderá ser solicitada uma ou mais das seguintes atividades.

\paragraph{Pré-teste do questionário (Componente~A).} Responder a um questionário sobre práticas agrícolas tradicionais em formato de entrevista assistida, com duração estimada de 30 a 45 minutos. Após responder, a equipe fará algumas perguntas sobre a clareza e a adequação das questões.

\paragraph{Entrevista (Componente~B).} Participar de uma conversa gravada em áudio, com duração estimada de 30 a 60 minutos, sobre suas práticas agrícolas, seus saberes tradicionais e o que você considera mais importante de ser protegido e documentado.

\paragraph{Aplicação do questionário completo (Componente~C).} Responder ao questionário adaptado sobre práticas agrícolas tradicionais, em formato de entrevista assistida, com duração estimada de 30 a 45 minutos.

\paragraph{Oficina devolutiva (Componente~C).} Participar de uma oficina em grupo, com duração estimada de 2 horas, na qual os resultados da pesquisa serão apresentados para avaliação e discussão pela comunidade.

\paragraph{Registro fotográfico.} Durante as atividades de campo, a equipe de pesquisa poderá realizar registros fotográficos de práticas agrícolas, sistemas de cultivo, paisagens e atividades coletivas, exclusivamente para fins de documentação científica. Nenhuma fotografia que permita sua identificação será publicada sem sua autorização expressa.

\section*{Quais são os riscos?}

Os riscos desta pesquisa são mínimos. É possível que algum tema abordado (perda de saberes, pressão sobre o território, dificuldades da comunidade) cause algum desconforto emocional. Caso isso aconteça, você poderá interromper sua participação a qualquer momento, sem precisar dar explicações e sem qualquer prejuízo. A equipe de pesquisa está preparada para acolhê-lo(a) e, se necessário, poderá encaminhá-lo(a) para serviços de apoio. Se você tiver dificuldade com leitura ou escrita, toda a aplicação será feita oralmente, sem nenhum constrangimento. Os dados coletados serão mantidos em sigilo e apresentados de forma que nenhum participante possa ser identificado. Os registros fotográficos serão tratados com o mesmo cuidado, sendo que imagens que permitam identificação individual somente serão utilizadas mediante sua autorização expressa.

\section*{Quais são os benefícios?}

Ao participar, você contribuirá para a documentação e proteção dos saberes agrícolas tradicionais da sua comunidade. Entre os benefícios previstos estão o acesso a um instrumento de documentação dos saberes em linguagem acessível, a participação ativa na definição do que é mais importante de ser protegido, o recebimento de um relatório com os resultados referentes à sua comunidade e o fortalecimento da base documental para processos junto a IPHAN, INPI e outros órgãos. Os mestres de saberes que participarem do comitê de adaptação e do painel de especialistas serão reconhecidos como coautores do instrumento.

\section*{Sigilo e proteção dos seus dados}

Suas respostas serão codificadas com um número, sem uso do seu nome. Os dados serão armazenados em ambiente digital seguro e criptografado, com acesso restrito à equipe de pesquisa, e serão mantidos por 5 anos após a publicação dos resultados, conforme a Resolução CNS nº~466/2012. Os áudios das entrevistas serão destruídos após transcrição e verificação. Os registros fotográficos serão armazenados em ambiente digital seguro e criptografado, com acesso restrito à equipe de pesquisa, sendo que imagens com identificação de participantes somente serão utilizadas mediante autorização expressa em termo próprio. Nenhuma informação que permita sua identificação será publicada. Os conhecimentos tradicionais específicos compartilhados serão tratados conforme a Lei nº~13.123/2015 (Marco Legal da Biodiversidade), assegurando a proteção contra apropriação indevida e a repartição justa de benefícios.

\section*{Participação voluntária e direito de desistência}

Sua participação é inteiramente voluntária. Você pode recusar-se a participar ou desistir a qualquer momento, sem necessidade de justificativa e sem qualquer prejuízo para você ou para a sua comunidade. A recusa em participar não afetará seu acesso a nenhum serviço ou programa.

\section*{Contatos}

Em caso de dúvidas sobre a pesquisa, entre em contato com o pesquisador responsável, Prof. Dr. Luiz Diego Vidal Santos, pelo telefone (\rule{1.5cm}{0.4pt}) \rule{3cm}{0.4pt} ou pelo e-mail \rule{5cm}{0.4pt}, ou com a doutoranda Catuxe Varjão de Santana Oliveira, pelo telefone (\rule{1.5cm}{0.4pt}) \rule{3cm}{0.4pt} ou pelo e-mail \rule{5cm}{0.4pt}.

Em caso de denúncia ou dúvida sobre questões éticas, entre em contato com o Comitê de Ética em Pesquisa da UFS (CEP/UFS), localizado na Cidade Universitária Prof. José Aloísio de Campos, Av. Marechal Rondon, s/n, Jardim Rosa Elze, São Cristóvão/SE, CEP 49100-000, pelo telefone (79) 3194-7208 ou pelo e-mail cepufs@academico.ufs.br.

\vspace{1cm}

\section*{Consentimento do(a) Participante}

Declaro que li (ou me foi lido) este Termo de Consentimento Livre e Esclarecido, que compreendi as informações nele contidas e que concordo voluntariamente em participar desta pesquisa. Recebi uma via deste documento e tive oportunidade de fazer perguntas, que foram respondidas de forma satisfatória.

\vspace{1.5cm}

\noindent
\begin{minipage}[t]{0.55\textwidth}
Local e data:\\[1.5cm]
\rule{8cm}{0.4pt}
\end{minipage}

\vspace{1cm}

\noindent
\begin{minipage}[t]{0.48\textwidth}
\centering
\rule{7cm}{0.4pt}\\
Nome e assinatura do(a) participante\\
(ou impressão digital)
\end{minipage}
\hfill
\begin{minipage}[t]{0.48\textwidth}
\centering
\rule{7cm}{0.4pt}\\
Nome e assinatura do(a) pesquisador(a)
\end{minipage}

\vspace{1.5cm}

\noindent
\textit{Em caso de participante não alfabetizado(a) ou com dificuldade de leitura, acrescentar assinatura de testemunha.}

\vspace{1cm}

\noindent
\begin{minipage}[t]{0.48\textwidth}
\centering
\rule{7cm}{0.4pt}\\
Nome e assinatura da testemunha\\
(quando aplicável)
\end{minipage}

\vspace{1cm}

\noindent
\footnotesize{Este documento é emitido em duas vias de igual teor, ficando uma com o(a) participante e outra com o pesquisador responsável.}

\end{document}
